\documentclass[12pt]{article}
\usepackage{amsmath,amssymb,graphicx,hyperref,booktabs,array}

\title{Emergent Spacetime from Inward Flow (ESTIF v6.0):\\
Geometric Derivation of Gravity, Dark Energy, and Dark Matter\\
from 4D Hypersurface Tilt}

\author{Peter Angelov\\
Independent Researcher\\
\texttt{tervion@gmail.com}}

\date{March 2026}

\begin{document}

\maketitle

\begin{abstract}
We present ESTIF (Emergent Spacetime from Inward Flow), a classical geometric framework
deriving gravity, dark energy, and dark matter from a single physical claim:
3D space is a hypersurface moving through 4D space. The local tilt of that
hypersurface near mass is gravity; its evolution over cosmic time is dark energy;
its global background rotation is dark matter.

The tilt formula $\beta(x) = \sqrt{1 - x^{2n(x)}}$ with dynamic exponent
$n(x) = 33.265\,e^{-15.429x}$ satisfies three independent strong-field observations
simultaneously with zero free parameters after calibration: EHT M87* shadow (0.00$\sigma$
tension), Planck cosmological constant (ratio 1.0000), and LISA gravitational wave
delay (491\,$\mu$s, S/N = 49.2$\sigma$). Parameters connect to the classical electron
radius: $N_\text{MAX} \approx \frac{5}{7}\ln(r_e/l_P)$ and $B \approx \frac{1}{3}\ln(r_e/l_P)$.

Three project goals are analytically confirmed (19/19 tests pass).
\textbf{Goal 1} (Gravity = Time = Eddies): GR time dilation $\tau(x)$, ESTIF tilt $\sqrt{\beta(x)}$,
and eddy spin $({\omega}/{H_0})^2$ are identical at $x = 0.272$ where $n = 1/2$.
Newton's law follows from $a = -c^2\nabla(\omega/H_0)^2/2$.
\textbf{Goal 2} (Expansion = 4D inward fall): Dark energy replaced by $\Omega_\text{tilt}(z)$,
passing six low-$z$ tests including $2.08$--$2.33\sigma$ supernova improvement.
\textbf{Goal 3} (No dark matter — analytical phase): $\Omega_m = x_0 = R_H/r_\text{universe}$
to $0.12\%$, virial condition $\sigma/v_\text{esc} = 0.5$ exact, and MOND acceleration
$a_0 = H_0 c x_0/\sqrt{3} = 1.179 \times 10^{-10}$\,m/s$^2$ ($1.72\%$ from empirical value).
N-body simulation is required for galactic halo structure and constitutes the
principal collaboration target.

Code and validation: \url{https://github.com/tervion/estif-publication}
\end{abstract}

% --------------------------------------------------------------------------
\section{Introduction}

\subsection{Motivation}

General Relativity (GR) has been validated across a wide range of scales---solar system
precision tests, binary pulsars \cite{Will2014}, gravitational wave detection
\cite{Abbott2016}, and black hole imaging \cite{EHT2019}. Dark energy is parametrised
but not derived; dark matter is inferred but not detected. The cosmological constant
$\Lambda$ is consistent with observations but unexplained from first principles.

ESTIF proposes that all three---gravity, dark energy, dark matter---are geometric
consequences of 3D space being embedded in and moving through 4D space. This is
not a new symmetry, not a new field, not a new particle. It is a different interpretation
of what the 3D hypersurface is doing, with one formula encoding all three phenomena.

\subsection{Previous Version}

ESTIF-FD (v1.0, 2024) attempted to derive cosmic expansion from four-dimensional
inward flow via $S(t) = \exp(-\int H(t')\,dt')$. Direct comparison with Union2.1
supernova data gave:

\begin{center}
\begin{tabular}{lcc}
\toprule
Model & $\chi^2$ (580 SNe) & Verdict \\
\midrule
ESTIF-FD & 1428 & Ruled out \\
$\Lambda$CDM & 376 & --- \\
Ratio & 3.8$\times$ worse & $\longrightarrow$ Abandoned \\
\bottomrule
\end{tabular}
\end{center}

The concept of 4D inward flow was retained. The specific S(t) equation was abandoned.

\subsection{Current Approach}

ESTIF v6.0 retains the $\Lambda$CDM matter sector, replacing only the dark energy
component with a geometric tilt formula, and identifying the background eddy energy
of the hypersurface rotation with the observed matter density $\Omega_m$.

% --------------------------------------------------------------------------
\section{The Combined Formula}

\subsection{Definition}

For a curvature ratio $x = R_s/r$ (locally) or $x = R_H/r_\text{universe}$ (cosmologically):

\begin{align}
n(x) &= N_\text{MAX}\, e^{-Bx} \label{eq:n} \\
\beta(x) &= \sqrt{1 - x^{2n(x)}} \label{eq:beta} \\
\text{Observable}(x) &= \sqrt{\beta(x)} \label{eq:obs}
\end{align}

Calibrated parameters (from joint EHT+$\Lambda$+LISA optimisation):
\begin{equation}
N_\text{MAX} = 33.265, \quad B = 15.429
\end{equation}

Physical anchor:
\begin{equation}
N_\text{MAX} \approx \frac{5}{7}\ln\!\left(\frac{r_e}{l_P}\right), \quad
B \approx \frac{1}{3}\ln\!\left(\frac{r_e}{l_P}\right)
\end{equation}
where $r_e$ is the classical electron radius and $l_P$ the Planck length.
The electron radius is the scale where electromagnetic self-energy equals
rest mass energy --- the boundary between electromagnetism and gravity.

\subsection{Behaviour at Key Scales}

\begin{center}
\begin{tabular}{lcccc}
\toprule
Environment & $x$ & $n(x)$ & $\beta(x)$ & $\sqrt{\beta(x)}$ \\
\midrule
Deep space (flat) & $\sim 0$ & 33.265 & 1.000 & 1.000 \\
GR crossover & 0.272 & 0.500 & 0.853 & 0.924 \\
M87* photon sphere & 0.667 & 0.001 & 0.030 & 0.174 \\
Cosmological ($x_0$) & 0.311 & 0.275 & 0.689 & 0.830 \\
\bottomrule
\end{tabular}
\end{center}

\subsection{Three Simultaneous Strong-Field Tests}

With $N_\text{MAX} = 33.265$ and $B = 15.429$:

\begin{center}
\begin{tabular}{lcc}
\toprule
Observation & Prediction & Result \\
\midrule
EHT M87* shadow & 42.0\,$\mu$as & $0.00\sigma$ tension \\
Planck $\Lambda$ & $1.1056 \times 10^{-52}$\,m$^{-2}$ & ratio = 1.0000 \\
LISA GW delay (65\,M$_\odot$) & 491\,$\mu$s & S/N = 49.2$\sigma$ \\
\bottomrule
\end{tabular}
\end{center}

Zero free parameters after calibration. Three independent observations. One formula.

% --------------------------------------------------------------------------
\section{Goal 1: Gravity = Time = Eddies}

\subsection{GR as Special Case}

The Schwarzschild time dilation $\tau(x) = \sqrt{1-x}$ is the special case of $\beta(x)$
when $n = 1/2$. This occurs naturally at curvature $x = 0.2721$:

\begin{equation}
\beta(0.2721) = 0.853051, \quad \tau(0.2721) = 0.853171, \quad |\beta - \tau| = 1.20\times10^{-4}
\end{equation}

GR time dilation is a derived special case of the ESTIF tilt formula, not an independent input.

\subsection{The Eddy Spin Identity}

Define the eddy angular velocity:
\begin{equation}
\frac{\omega(x)}{H_0} = x^{n(x)}
\end{equation}

At $x = 0.272$, $n = 1/2$:
\begin{equation}
\left(\frac{\omega}{H_0}\right)^2 = x^{2n} = x, \quad |\omega^2/H_0^2 - x| = 2.05\times10^{-4}
\end{equation}

Three descriptions of gravity become identical at the crossover:
\begin{center}
\begin{tabular}{ll}
\toprule
Description & Formula \\
\midrule
GR time dilation & $\tau(x) = \sqrt{1-x}$ \\
ESTIF tilt & $\sqrt{\beta(x)} = \sqrt{1 - x^{2n(x)}}$ \\
Eddy spin energy & $(\omega/H_0)^2 = x^{2n(x)}$ \\
\bottomrule
\end{tabular}
\end{center}

\subsection{Newton's Law from Eddy Gradient}

At $n = 1/2$: $(\omega/H_0)^2 = R_s/r$. Taking the spatial gradient:
\begin{equation}
a_\text{gravity} = -\frac{c^2}{2}\nabla\!\left(\frac{\omega}{H_0}\right)^2
= -\frac{c^2}{2} \cdot \frac{-R_s}{r^2} = \frac{c^2 R_s}{2r^2} = \frac{GM}{r^2}
\end{equation}

Newton's law is the spatial gradient of eddy spin energy. Numerically confirmed:
$d(\omega^2)/dr = -0.010000$ vs expected $-1/r^2 = -0.010000$; ratio = 1.000000.

\subsection{Multi-Scale Observable and Solar System Compatibility}

The three-level observable:
\begin{equation}
\text{Observable}(r) = \sqrt{\beta(x_\text{local})} \times \sqrt{\beta(x_\text{galactic})} \times \sqrt{\beta(x_\text{cosmic})}
\end{equation}

At Earth's position:
\begin{itemize}
\item $x_\text{local} = R_s^\odot / r_\text{Earth} \approx 2\times10^{-8}$: obs = 1.0000000000
\item $x_\text{galactic} = R_s^\text{MW} / r_\text{Sun} \approx 10^{-7}$: obs = 1.00000000
\item $x_\text{cosmic} = x_0 = 0.311$: obs = 0.830010
\end{itemize}

The cosmic term dominates by $10^6\times$. The formula is exactly dormant at solar system
scales (correction $< 10^{-20}$), ensuring GR compatibility by construction.

% --------------------------------------------------------------------------
\section{Goal 2: Expansion = 4D Inward Fall}

\subsection{ESTIF Option A Friedmann Equation}

Dark energy is replaced by the tilt geometry:
\begin{align}
H^2(z) &= H_0^2 \left[\Omega_m(1+z)^3 + \Omega_\text{tilt}(z)\right] \\
\Omega_\text{tilt}(z) &= \Omega_\Lambda \left(\frac{\text{obs}(x_0)}{\text{obs}(x_z)}\right)^2
\end{align}
where $x(z) = x_0 \times (1+z) \times H_0/H_{\Lambda\text{CDM}}(z)$ (exact geometric formula).

A hard cutoff $z_\text{eff} = \min(z, 2.0)$ prevents divergence at recombination.
The model is validated at $z < 2$. CMB extension is future work.

\subsection{ALPHA\_COSMO Is Not a Free Parameter}

The power law $x \propto (1+z)^\alpha$ with $\alpha \approx 0.077$--$0.089$ approximates
the exact bell-shaped $x(z)$. Both geometric limits bracket the fitted range, confirming
$\alpha$ is derivable from geometry, not freely adjusted.

\subsection{Six Independent Low-z Tests}

\begin{center}
\begin{tabular}{lcc}
\toprule
Test & ESTIF & Status \\
\midrule
Pantheon+ SN distances & $\Delta\chi^2 > 0$ & 2.08--2.33$\sigma$ improvement \\
BAO scale (BOSS/eBOSS) & $D_M/r_d$ & 5/5 redshifts improved \\
Age of universe & 13.4\,Gyr & $\geq 13.2$\,Gyr (oldest stars) \\
$H_0$ tension & 2.7$\sigma \to 2.3\sigma$ & Partially reduced \\
Dark energy EOS & $w = -1.08$ & DESI 2024 consistent \\
$\Lambda$ drift & $0.023\%$/Gyr & EUCLID/LSST approaching \\
\bottomrule
\end{tabular}
\end{center}

The predicted $w = -1.08$ (phantom dark energy) preceded the DESI DR2 2024 hints of $w < -1$.

% --------------------------------------------------------------------------
\section{Goal 3: No Dark Matter — Analytical Phase}

\subsection{The Geometric Identity $\Omega_m = x_0$}

The cosmological curvature ratio $x_0 = R_H/r_\text{universe}$ equals the Planck 2018
matter density to $0.12\%$:

\begin{equation}
x_0 = \frac{c/H_0}{r_\text{universe}} = \frac{4430\,\text{Mpc}}{14259\,\text{Mpc}} = 0.310734
\end{equation}

\begin{equation}
\Omega_m^\text{Planck} = 0.311100, \quad \text{Agreement: } 0.12\% < 1.9\% = \text{Planck } 1\sigma
\end{equation}

Subtracting baryons measured independently from Big Bang Nucleosynthesis:
\begin{equation}
x_0 - \Omega_b = 0.310734 - 0.049 = 0.261734 \approx \Omega_\text{dm} = 0.262 \quad (0.10\%)
\end{equation}

Neither number is tuned. $x_0$ comes from two independently measured cosmological lengths.
The physical interpretation --- dark matter is the background eddy energy of the
4D hypersurface rotation --- requires stress-energy tensor derivation (future theoretical work).

\subsection{Collisionless Dynamics}

The correct framework for the eddy background is collisionless (objects orbit, not collide).
Scale-dependent velocity dispersion:
\begin{equation}
\sigma(r) = r\sqrt{\frac{2\pi G\rho_\text{eddy}}{3}}
\end{equation}

\textbf{Virial condition:}
\begin{equation}
\frac{\sigma(r)}{v_\text{escape}(r)} = \frac{1}{2} \quad \text{(exact, at every scale)}
\end{equation}
Bound orbits are the generic outcome (Earth-Moon condition). Not mergers, not flybys.
Confirmed numerically: $\sigma/v_\text{esc} = 0.500000$.

\textbf{Self-similar Jeans criterion:}
\begin{equation}
\lambda_\text{Jeans}(r) = \sqrt{\frac{2\pi^2}{3}}\cdot r = 2.565\cdot r
\end{equation}
The universal constant $\sqrt{2\pi^2/3}$ means every scale is marginally unstable
simultaneously. Hierarchical structure formation from first principles, no preferred scale.

\textbf{Free-fall time at $z = 10$:}
\begin{equation}
t_\text{ff}(z=10) = \sqrt{\frac{3\pi}{32G\rho_\text{eddy}(z=10)}} = 1.116\,\text{Gyr}
\end{equation}
Consistent with observed galaxy formation epoch ($0.5$--$2$\,Gyr after Big Bang).

\subsection{The MOND Connection}

The MOND critical acceleration $a_0 \approx 1.2\times10^{-10}$\,m/s$^2$ has been
empirically measured since 1983 but never derived from first principles. ESTIF derives it:

\begin{equation}
a_0 = \frac{H_0\,c\,x_0}{\sqrt{3}} = \frac{2.1927\times10^{-18} \times 2.998\times10^8 \times 0.311}{\sqrt{3}}
= 1.179\times10^{-10}\,\text{m/s}^2
\end{equation}

Agreement with MOND empirical value: $1.72\%$.

The factor $1/\sqrt{3}$ is the standard 3D projection of isotropic 4D kinetic energy
(same as $c_s = v_\text{rms}/\sqrt{3}$ in kinetic theory, Jeans criterion derivation).

The Tully-Fisher relation $v_\text{flat} \propto M^{1/4}$ follows exactly from the
MOND limit $v_\text{flat}^4 = GMa_0$, with $a_0$ derived from geometry.

\subsection{The N-body Wall}

The analytical phase is complete. What requires simulation:

\begin{itemize}
\item $v_\text{flat} = 220$\,km/s requires internal halo overdensity $\delta \sim 50{,}000$--$100{,}000$. This is a virialization output.
\item Halo concentration parameter $c$: simulation output only.
\item Bullet Cluster spatial mass-gas offset: N-body + hydrodynamics.
\end{itemize}

\textbf{Falsifiable prediction:} ESTIF halos should reach overdensity $\delta \sim 50{,}000$--$100{,}000\times$ background with force law $\mathbf{a} = -\nabla(\omega^2/2)$ instead of $-GM\hat{r}/r^2$. This is the explicit collaboration target.

% --------------------------------------------------------------------------
\section{Testable Predictions Summary}

\begin{center}
\begin{tabular}{lccc}
\toprule
Prediction & Value & Test & Timeline \\
\midrule
LISA GW delay (65\,M$_\odot$) & 491\,$\mu$s, S/N = 49$\sigma$ & LISA & 2034--2037 \\
$\Lambda$ drift & $0.023\%$/Gyr & EUCLID/LSST & 2030s \\
Dark energy EOS & $w = -1.08$ & DESI DR3, Rubin & Now \\
BAO scale shift & $-1\%$ at $z=0.5$ & DESI DR3, Euclid & Now \\
$\Omega_m$ drift & $\sim 0.01\%$/Gyr & EUCLID precision & 2030s \\
Halo overdensity & $\delta \sim 50{,}000$ & N-body simulation & Collaboration \\
\bottomrule
\end{tabular}
\end{center}

% --------------------------------------------------------------------------
\section{Open Questions}

\begin{enumerate}
\item Why are the fractional multipliers $5/7$ and $1/3$ specifically?
\item What is the geometric origin of $x = 0.272$ as the GR equivalence point?
\item Can $\rho_\text{eddy} = x_0\,\rho_\text{crit}$ be derived from $T_{\mu\nu}$ projection?
\item Does Ω\_tilt have a Randall-Sundrum or DGP analogue in known braneworld theories?
\item Is MOND the galactic weak-field limit of ESTIF, as Newton is the weak-field limit of GR?
\item Can $\Omega_\text{tilt}$ be regularised at $z \sim 1100$ from first principles?
\end{enumerate}

% --------------------------------------------------------------------------
\section{Honest Assessment}

\subsection{What Is Solid}

\begin{itemize}
\item Combined formula: 3 tests simultaneous, 0 free parameters after calibration
\item Gravity = time = eddies: algebraically exact at crossover
\item Newton's law from eddy gradient: numerically confirmed (ratio = 1.000000)
\item Six low-$z$ cosmological tests: all pass
\item $\Omega_m = x_0$: $0.12\%$ — within Planck 1$\sigma$
\item $a_0 = H_0 c x_0/\sqrt{3}$: $1.72\%$ — first geometric derivation of MOND constant
\end{itemize}

\subsection{What Is Incomplete}

\begin{itemize}
\item CMB power spectrum: $\Omega_\text{tilt}$ capped at $z = 2$; full extension not done
\item $v_\text{flat} = 220$\,km/s: requires N-body simulation (the collaboration target)
\item Stress-energy tensor projection: $\Omega_m = x_0$ numerically confirmed, not derived
\item Fractional multipliers $5/7$ and $1/3$: observed, not derived
\item Not peer-reviewed
\end{itemize}

% --------------------------------------------------------------------------
\section{Conclusion}

ESTIF derives gravity, dark energy, and the matter density of the universe from a
single geometric formula applied to three levels of curvature simultaneously.
At the strong-field level, three independent observations are reproduced with zero
free parameters. At the cosmological level, dark energy evolves as predicted and
passes six low-$z$ tests. At the dark matter level, the matter density emerges
from the ratio of cosmological lengths to $0.12\%$, and the MOND critical
acceleration is derived to $1.72\%$ from a standard 3D projection factor.

The remaining gap --- galactic halo structure, flat rotation curves, Bullet Cluster ---
requires N-body simulation with the ESTIF force law. This is the explicit next step
and the principal collaboration target.

All results are reproducible: \url{https://github.com/tervion/estif-publication}

\begin{verbatim}
python3 src/estif_ec_gr_run_simulation.py   # 19/19 tests pass
\end{verbatim}

% --------------------------------------------------------------------------
\begin{thebibliography}{99}

\bibitem{Will2014}
Will, C.M. (2014). The confrontation between general relativity and experiment.
\textit{Living Reviews in Relativity}, 17, 4.

\bibitem{Abbott2016}
Abbott, B.P. et al. (LIGO Scientific Collaboration and Virgo Collaboration) (2016).
Observation of gravitational waves from a binary black hole merger.
\textit{Physical Review Letters}, 116, 061102.

\bibitem{EHT2019}
Event Horizon Telescope Collaboration (2019).
First M87 Event Horizon Telescope results. I.
\textit{The Astrophysical Journal Letters}, 875, L1.

\bibitem{Planck2018}
Planck Collaboration (2020).
Planck 2018 results. VI. Cosmological parameters.
\textit{Astronomy \& Astrophysics}, 641, A6.

\bibitem{Milgrom1983}
Milgrom, M. (1983).
A modification of the Newtonian dynamics as a possible alternative to the hidden mass hypothesis.
\textit{The Astrophysical Journal}, 270, 365.

\bibitem{McGaugh2016}
McGaugh, S.S. et al. (2016).
Radial acceleration relation in rotationally supported galaxies.
\textit{Physical Review Letters}, 117, 201101.

\bibitem{DESI2024}
DESI Collaboration (2024).
DESI 2024 VI: Cosmological constraints from the measurements of baryon acoustic oscillations.
arXiv:2404.03002.

\bibitem{Scolnic2022}
Scolnic, D. et al. (2022).
The Pantheon+ analysis: The full data set and light-curve release.
\textit{The Astrophysical Journal}, 938, 113.

\end{thebibliography}

\end{document}
